\addcontentsline{toc}{section}{Введение}

\section*{Введение}

% База данных — совокупность данных, хранимых как в соответствии со схемой данных так и в произвольном порядке, манипулирование которыми выполняют в соответствии с правилами средств моделирования данных.
% Система управления базами данных — совокупность программных и лингвистических средств общего или специального назначения, обеспечивающих управление созданием и использованием баз данных.

% Чтобы создать базу данных, сначала проектируют её структуру: определяют,
% какие данные хранить и как они связаны между собой. Затем схему переносят
% в СУБД и заполняют данными. Если база спроектирована правильно, данные
% не дублируются, запросы выполняются быстро и всё остаётся в порядке
% даже когда с базой работают одновременно несколько человек.

% Данная работа посвящена проектированию и реализации реляционной базы
% данных для предметной области «Геомониторинг автопарка каршеринга».
% Современные сервисы каршеринга управляют сотнями автомобилей, и отслеживание их местонахождения, контроль парковочных сессий,
% а также своевременная реакция на тревожные события — угоны, выезды
% за разрешённые зоны, потерю связи — требуют надёжной системы хранения
% и быстрого доступа к данным.

% В качестве СУБД используется PostgreSQL~14; заполнение базы данных
% реализовано на языке Go с использованием библиотек \textbf{pgx/v5}
% и \textbf{golang-migrate}.
